\section{Experimental design in full}
\label{app:design}

\paragraph{Model.} All games in E0 to E2, E4, the grounded runs and the two benchmarks use \texttt{gpt-6-luna} on both sides with reasoning off (temperature 0.7, top-$p$ 0.8, at most 700 output tokens per move); the extractor and the judge use the same model with low reasoning. E3 adds \texttt{claude-haiku-5-5}, \texttt{claude-sonnet-5-5} and \texttt{gpt-6.1-sol}, served through a university proxy that accepts only the default temperature, with thinking and reasoning disabled (hidden reasoning otherwise consumes the reply budget), and \texttt{gpt-6-luna} at its default temperature as the control for the sampling difference; \texttt{claude-sonnet-5-5} needed a 1{,}500-token cap under open information, where its replies run long. The judge labeled every E1 game and, after the fact, E2's base condition, E4, the grounded run and the E3 arms (a style per agent from the public transcript), as the secondary measure. Instances are drawn from seeds that are disjoint between the development set and every test set, and each configuration, with its predictions, was frozen before its run.

\paragraph{E0: manipulation check.} On a development set of 10 instances in the core regime, Developer 1 receives each strategy in each operationalization against a \neutral{} partner (12 cells, 120 games per round). The rules were revised over three rounds on this set until compliance was high for every strategy (\cref{sec:e0}); the test sets use other seeds.

\paragraph{E1: the strategy matrix.} All 36 ordered pairs of \neutral, \competing, \collaborating, \compromising, \accommodating{} and \avoiding{} as rules, the five prompted sentences on the diagonal and no guidance on both sides, on 40 instances with the first mover alternating (42 cells, 1{,}680 games). Primary analyses, registered before the run: joint payoff by pair against the \neutral{} pair and the own-payoff matrix, both paired by instance with exact sign tests; secondary: best responses per column, the attractor of replicator dynamics on the empirical $6\times6$ game, and the exploitability of each strategy. Intention to treat is primary (every game counts whether or not the agent complied) and compliant games are a secondary cut.

\paragraph{E2: the condition map.} Ten conditions, each one knob moved from the E1 regime (\cref{tab:conditions} in \cref{app:conditions}): urgency ($T{=}4$), a cost of 3 points per move, a reply budget of 250 tokens, nothing to trade ($a{=}0$), everything to trade ($a{=}1$), every issue contested ($p{=}1$), unequal stakes (1 against 0.3), unequal power (outside options 45 against 15), open information (scores visible), and the base regime again. Each condition plays the 24 cells that read the map (the six diagonal pairs and every strategy against \neutral{} and against \competing{} in both seats) on 20 shared instances: 480 games per condition, 4{,}800 in all. For each condition we report the best strategy in three readings, symmetric play, adopted alone against the \neutral{} partner and against a \competing{} partner, by points and by net payoff over $\lam$, with seats pooled.

\paragraph{E4: framing.} E1's 40 instances under the three skins, with the rule diagonal and the two unilateral deviations against \neutral{} (8 cells, 960 games), so that every game has an abstract twin in E1 played with byte-identical rule texts.

\paragraph{Grounded.} E2's 24 cells on the 85 ParallACT-derived instances (2{,}040 games), the same structures replayed under the abstract skin as a control (2{,}040 games), a replay of the \collaborating{} cells with the single-contested-issue clause of the rule (680 games per rendering), and the instances with the ticket context on a second model, \texttt{claude-haiku-5-5}, with the 16 cells that read the map without a \competing{} opponent (1{,}360 games).

\paragraph{E3: models.} E1's eight reading cells (the rule diagonal, \competing{} and \collaborating{} against \neutral) on E1's 40 instances with byte-identical rule texts, on each of the four arms above (1{,}280 games), so that every game has an E1 twin; then the two conditions that moved the map most on \texttt{gpt-6-luna}, open information and unequal stakes, with all 24 E2 cells on E2's 20 instances for \texttt{claude-sonnet-5-5} and \texttt{gpt-6.1-sol} (1{,}920 games). \textbf{Cross-model pairs.} A \texttt{gpt-6-luna} agent against a \texttt{claude-sonnet-5-5} agent on E1's 40 instances, with the 16 cells that read the map (the diagonal and every strategy against \neutral{} in both seats) and the seats' models swapped in a second arm that plays the mixed \neutral{} pair and \competing{} and \collaborating{} against \neutral{} (840 games), so that each model adopts each strategy against the other's default agent. The question is whether what a rule earns its user depends on the user's model or on the partner's.

\paragraph{External benchmarks.} Two published multi-issue negotiation games reduced to the environment's instance format, so that the same rules, checker and measures run unchanged. CaSiNo \citep{chawla2021casino}: two campers split three packages each of food, water and firewood, worth 5, 4 and 3 points a package by priority, with a walk-away worth 5 of 36; all 36 ordered pairs of priority orders are the instances, each item a decision with four options. LLM-Deliberation \citep{abdelnabi2024cooperation}: six parties negotiate five issues with three to five options, each with a score table and an acceptance threshold; we keep every pair of parties with a contested decision in its four games (57 pairs), their score tables over all five issues, an outside option of 30 percent of each party's maximum as in the core, and the game's own scenario text. Both play the 16 reading cells on \texttt{gpt-6-luna} with E1's settings (576 and 912 games), so that every difference from the core is the game's, not the model's. Neither reduction uses dialogue text.

\paragraph{Cost of the study.} About 19{,}800 games for \$78 in model calls. The \texttt{gpt-6-luna} studies cost under \$10 together (E1 \$1.51 including the judge, E2 \$2.50, the grounded runs \$2.33, E4 \$0.56, E0 \$0.33, the two benchmarks \$0.79, the post-hoc labels \$1.90); the other three models cost \$69 (\texttt{claude-sonnet-5-5} and \texttt{gpt-6.1-sol} \$0.015 to \$0.03 a game, \texttt{claude-haiku-5-5} \$0.001).

